\scanpage{98}
liés sss $(s,t)\in\mathrm{Im}\,\varphi_{\{\alpha,\beta\}}(\widetilde\Gamma)$.

Étant donné des données a) b) c), satisfaisant la condition $(\alpha)$, le graphe qu'on vient de décrire satisfait les conditions \textcircled{1} (Axiome 1) et \textcircled{2} (L'axiome 2 pour le \inserted{seul} triangle $t_0$). L'ens \inserted{$T$} des triangles de $(S,A)$ est alors donné, mod bijection, par
\[
T\simeq\underbrace{\{t_0\}}_{\substack{\text{triangle}\\\text{de base}}}\amalg
\underbrace{\coprod_{\alpha\in t_0}E_\alpha}_{\substack{\text{triangles}\\\text{verticaux}}}\amalg
\underbrace{\widetilde\Gamma}_{\substack{\text{triangles (horiz.)}\\\text{supérieurs}}}
\]
où on pose \deleted{\unsure{$E$}}
\[
E_\alpha=\widetilde E_\alpha/\sigma_\alpha
\]
(de sorte que $\widetilde E_\alpha$ est un rev. de degré 2 de $E_\alpha$). Pour $\alpha\in t_0$, un élément $u$ de $E_\alpha$, i.e. $u\in\unsure{\mathcal P|\mathfrak P}_2(\widetilde E_\alpha)$ formé d'une orbite sous $\sigma_\alpha$, \unsure{son|le} triangle \inserted{(vertical)} associé est par définition \deleted{\unsure{$\{$}}$u\unsure{\amalg|\cup}\{\alpha\}$. Pour $\deleted{\unsure{\gamma\in\widetilde\Gamma}}\tau\in\widetilde\Gamma$, le triangle associé est par définition l'ens \deleted{\unsure{des}} $\{\varphi_\alpha(\tau)\}_{\alpha\in t_0}$.

\inserted{Rien n'empêche pour l'instant qu'on ait $\widetilde\Gamma=\varnothing$ ! Mais} On se propose maintenant d'exprimer l'axiome 2 pour \uline{ts} les triangles de $S$ (pas seulement $t_0$). Prenons d'abord les
% archive sheet number: 97
